%%%%%%%%%%%%%%%%%%%%%%%%%%%%%%%%%
%
%       EXERCÍCIO
%
%%%%%%%%%%%%%%%%%%%%%%%%%%%%%%%%%

\ifdefstring{\atividade}{prova}{%
    \ifdefstring{\modo}{objetivo}{%
        \renewcommand{\valorquestao}{\ValorQObj\ ponto}
    }{%
        \renewcommand{\valorquestao}{\ValorQDisc\ pontos}
    }%
}%

\begin{exercicioBanco}[\valorquestao]
Em uma festa, $35$ pessoas discutiam sobre dois filmes, $A$ e $B$. 
Sabe-se que:
%
\begin{enumerate}[\(\bullet\)]
    \begin{multicols}{2}
        \item $15$ pessoas assistiram ao filme $A$;
        \item $8$ pessoas assistiram aos dois filmes;
        \item $10$ pessoas não assistiram a nenhum dos dois filmes.
    \end{multicols}
\end{enumerate}
%
Sabendo que todas as $35$ pessoas opinaram, o número de pessoas que assistiram ao filme $B$ é:


% Define as alternativas
\newcommand{\alternativas}{%
    \begin{center}
        \begin{tabularx}{\textwidth}{XXXXX}
            (a) \(7\). &
            (b) \(8\). &
            (c) \(10\). &
            (d) \(18\). &
            (e) \(25\).
        \end{tabularx}
    \end{center}
}

% Define a resposta correta
\newcommand{\resposta}{D}

% Lógica condicional para exibição
\ifdefstring{\atividade}{lista}{%
    \alternativas
    \vspace{0.5em}
    
    \noindent\textbf{Resposta:} letra \textbf{\resposta}.
}{%
    \ifdefstring{\modo}{objetivo}{%
        \alternativas
    }{%
        % modo = discursiva → não mostra alternativas
    }
}
\end{exercicioBanco}

